\documentclass[12pt, a4paper]{article}
\usepackage[UTF8]{ctex}

\usepackage{amsmath, amssymb, amsthm}
\usepackage{geometry}
\usepackage{fancyhdr}
\usepackage{enumitem}
\usepackage{xcolor}
\usepackage{hyperref}

\geometry{margin=1in}
\pagestyle{fancy}
\fancyhf{}
\rhead{线性代数（A）}
\lhead{李生丹}
\rfoot{第 \thepage 页}

\newcommand{\problem}[1]{\subsection*{题 #1}}
\newcommand{\solution}{\paragraph{解.}}

\definecolor{solutioncolor}{RGB}{0, 100, 0}

\begin{document}

\begin{center}
    {\LARGE\bfseries 线性代数 第 3 次作业 - 自动求解} \\[0.5em]
    {\large 线性代数（A）} \\[0.3em]
    李生丹 \quad | \quad 2026-10-05
\end{center}

\hrule
\vspace{1em}
\problem{1}

计算下列矩阵的行列式：  
2 1 3
 
 
A = 4 5 6
7 8 9

\solution

矩阵为 $A=\begin{pmatrix}2&1&3\\4&5&6\\7&8&9\end{pmatrix}$。

按第一行展开：$\det A=2\begin{vmatrix}5&6\\8&9\end{vmatrix}-1\begin{vmatrix}4&6\\7&9\end{vmatrix}+3\begin{vmatrix}4&5\\7&8\end{vmatrix}$。

计算二阶行列式：$\begin{vmatrix}5&6\\8&9\end{vmatrix}=45-48=-3$，$\begin{vmatrix}4&6\\7&9\end{vmatrix}=36-42=-6$，$\begin{vmatrix}4&5\\7&8\end{vmatrix}=32-35=-3$。

代入得 $\det A=2(-3)-1(-6)+3(-3)=-6+6-9=-9$。

\textbf{Answer:}
\[\boxed{-9}\]

\vspace{1em}
\hrule
\vspace{1em}

\problem{2}

求下列矩阵的逆矩阵： ( )
2 1
A =
5 3

\solution

矩阵 $A=\begin{pmatrix}2&1\\5&3\end{pmatrix}$，先求行列式 $\det A=2\cdot3-1\cdot5=1$。

因为 $\det A\neq0$，逆矩阵存在。

二阶矩阵逆公式：$A^{-1}=\frac{1}{\det A}\begin{pmatrix}d&-b\\-c&a\end{pmatrix}$。

代入 $a=2,b=1,c=5,d=3$，得 $A^{-1}=\begin{pmatrix}3&-1\\-5&2\end{pmatrix}$。

\textbf{Answer:}
\[\boxed{\begin{pmatrix}3&-1\\-5&2\end{pmatrix}}\]

\vspace{1em}
\hrule
\vspace{1em}

\problem{3}

求矩阵的特征值与特征向量：
( )
4 1
A =
2 3

\solution

矩阵 $A=\begin{pmatrix}4&1\\2&3\end{pmatrix}$，特征方程为 $\det(A-\lambda I)=0$。

$\det\begin{pmatrix}4-\lambda&1\\2&3-\lambda\end{pmatrix}=(4-\lambda)(3-\lambda)-2=\lambda^2-7\lambda+10=0$。

解得特征值 $\lambda_1=5,\lambda_2=2$。

对 $\lambda_1=5$：解 $(A-5I)x=0$，即 $\begin{pmatrix}-1&1\\2&-2\end{pmatrix}x=0$，得 $x_1=x_2$，取特征向量 $\begin{pmatrix}1\\1\end{pmatrix}$。

对 $\lambda_2=2$：解 $(A-2I)x=0$，即 $\begin{pmatrix}2&1\\2&1\end{pmatrix}x=0$，得 $2x_1+x_2=0$，取特征向量 $\begin{pmatrix}1\\-2\end{pmatrix}$。

\textbf{Answer:}
\[\boxed{\lambda_1=5,\ v_1=\begin{pmatrix}1\\1\end{pmatrix};\quad \lambda_2=2,\ v_2=\begin{pmatrix}1\\-2\end{pmatrix}}\]

\vspace{1em}
\hrule
\vspace{1em}

\problem{4}

解线性方程组：  
1 2 3 6
 
 
 2 4 6 12 
1 1 1 3

\solution

增广矩阵为 $\begin{pmatrix}1&2&3&6\\2&4&6&12\\1&1&1&3\end{pmatrix}$。

第二行减去第一行的 $2$ 倍：$R_2-2R_1=(0,0,0,0)$，说明第二个方程与第一个方程等价。

第三行减去第一行：$R_3-R_1=(0,-1,-2,-3)$，即 $-y-2z=-3$，得 $y=3-2z$。

由第一行 $x+2y+3z=6$，代入 $y=3-2z$：$x+2(3-2z)+3z=6$，得 $x+6-4z+3z=6$，即 $x=z$。

令 $z=t$，则 $x=t,\ y=3-2t,\ z=t$，$t$ 为任意实数。

\textbf{Answer:}
\[\boxed{\begin{cases}x=t\\y=3-2t\\z=t\end{cases},\ t\in\mathbb{R}}\]

\vspace{1em}
\hrule
\vspace{1em}

\problem{5}

求下列矩阵的秩：  
1 2 3
 
 
A = 2 4 6
3 6 9

\solution

矩阵 $A=\begin{pmatrix}1&2&3\\2&4&6\\3&6&9\end{pmatrix}$。

第二行是第一行的 $2$ 倍，第三行是第一行的 $3$ 倍，因此所有行都成比例。

对矩阵作初等行变换：$R_2-2R_1=0$，$R_3-3R_1=0$，只剩第一行非零。

非零行数为 $1$，故秩为 $1$。

\textbf{Answer:}
\[\boxed{1}\]

\vspace{1em}
\hrule
\vspace{1em}

\problem{6}

已知矩阵 ( ) ( )
1 2 2 0
A = , B =
3 4 1 3
计算矩阵乘积 AB。

\solution

$A=\begin{pmatrix}1&2\\3&4\end{pmatrix}$，$B=\begin{pmatrix}2&0\\1&3\end{pmatrix}$。

$(AB)_{11}=1\cdot2+2\cdot1=4$，$(AB)_{12}=1\cdot0+2\cdot3=6$。

$(AB)_{21}=3\cdot2+4\cdot1=10$，$(AB)_{22}=3\cdot0+4\cdot3=12$。

所以 $AB=\begin{pmatrix}4&6\\10&12\end{pmatrix}$。

\textbf{Answer:}
\[\boxed{\begin{pmatrix}4&6\\10&12\end{pmatrix}}\]

\vspace{1em}
\hrule
\vspace{1em}

\problem{7}

用 Gram-Schmidt 方法将下列向量组正交化：
   
1 1
   
   
α
1
= 1, α
2
= 0
0 1

\solution

给定 $\alpha_1=\begin{pmatrix}1\\1\\0\end{pmatrix}$，$\alpha_2=\begin{pmatrix}1\\0\\1\end{pmatrix}$。

取 $\beta_1=\alpha_1=\begin{pmatrix}1\\1\\0\end{pmatrix}$。

计算投影系数：$\frac{\alpha_2\cdot\beta_1}{\beta_1\cdot\beta_1}=\frac{1\cdot1+0\cdot1+1\cdot0}{1^2+1^2+0^2}=\frac{1}{2}$。

$\beta_2=\alpha_2-\frac{1}{2}\beta_1=\begin{pmatrix}1\\0\\1\end{pmatrix}-\frac{1}{2}\begin{pmatrix}1\\1\\0\end{pmatrix}=\begin{pmatrix}1/2\\-1/2\\1\end{pmatrix}$。

可取等价的正交向量组 $\beta_1=\begin{pmatrix}1\\1\\0\end{pmatrix}$，$\beta_2=\begin{pmatrix}1\\-1\\2\end{pmatrix}$（后者为前者的 $2$ 倍）。

\textbf{Answer:}
\[\boxed{\beta_1=\begin{pmatrix}1\\1\\0\end{pmatrix},\quad \beta_2=\begin{pmatrix}1\\-1\\2\end{pmatrix}}\]

\vspace{1em}
\hrule
\vspace{1em}

\problem{8}

判断矩阵 A 是否可对角化；若可以，求出对角化所需的矩阵 P 与对角矩阵 D：
( )
4 1
A =
2 3

\solution

矩阵 $A=\begin{pmatrix}4&1\\2&3\end{pmatrix}$，特征方程为 $\det(A-\lambda I)=\lambda^2-7\lambda+10=0$。

特征值为 $\lambda_1=5,\lambda_2=2$，两个互异特征值，故 $A$ 可对角化。

对 $\lambda_1=5$：$(A-5I)x=0$ 得 $x_1=x_2$，取特征向量 $p_1=\begin{pmatrix}1\\1\end{pmatrix}$。

对 $\lambda_2=2$：$(A-2I)x=0$ 得 $2x_1+x_2=0$，取特征向量 $p_2=\begin{pmatrix}1\\-2\end{pmatrix}$。

令 $P=\begin{pmatrix}1&1\\1&-2\end{pmatrix}$，$D=\begin{pmatrix}5&0\\0&2\end{pmatrix}$，则 $A=PDP^{-1}$。

\textbf{Answer:}
\[\boxed{P=\begin{pmatrix}1&1\\1&-2\end{pmatrix},\quad D=\begin{pmatrix}5&0\\0&2\end{pmatrix}}\]

\vspace{1em}
\hrule
\vspace{1em}


\vspace{2em}
\noindent\textit{Auto-generated: 8/8 problems solved by SymPy.}
\end{document}
